% 来源：OEBPS/Text/01_21.xhtml（OPF spine 顺序；仅第一册）
\cleardoublepage
\BookEntry[贾雷德·戴蒙德《枪炮、病菌与钢铁》]{贾雷德·戴蒙德《枪炮、病菌与钢铁》}

\BookCover{images/书单08枪炮、病菌与钢铁.jpg}{[美] 贾雷德·戴蒙德著| 谢延光译| 上海译文出版社2014年版}

\begin{BookSummaryBox}
{\centering \textbf{理解人类社会发展史的重要著作\\  1988年美国普利策奖获奖图书}\par}
\end{BookSummaryBox}

\subsubsection*{这本书谈什么？}

作为一本关于人类发展史的书，它讨论了一个很诱人的问题：为什么在近代把美洲人、亚洲人、澳洲人挨个入侵一遍的是欧洲人，而不是澳大利亚土著，不是印第安人，不是新几内亚人，也不是非洲人，甚至也不是在工业革命前在世界占主导地位的中国呢？从这个问题开始，作者试着从人类走出非洲开始，从走向每片大陆的时间切入，把农业的兴起、大型食草哺乳动物的驯服、发明和技术的传播速度作为主要讨论对象，对这个问题进行了阐释。
\par

贾雷德·戴蒙德作为演化生物学家，解释了事实上有助于形成历史最广泛模式的环境因素，以震撼人心的力量摧毁了以种族主义为基础的人类史理论。
\par

\subsubsection*{芒格为什么谈起这本书？}

芒格和巴菲特在进行商业分析和评估时，会采用“多元思维模型”，他们认为：
\par

\begin{quotation}
你必须知道重要学科的重要理论，并经常使用它们——要全部都用上，而不是只用几种。大多数人都只使用学过的一个学科的思维模型，比如说经济学，试图用一种方法来解决所有问题。你知道谚语是怎么说的：“在手里拿着铁锤的人看来，世界就像一颗钉子。”这是处理问题的一种笨办法。
\par
\end{quotation}

然而，英雄所见略同。芒格和巴菲特并不是仅有的两个使用非商业模型取得巨大成功的精英投资者。传奇式的债券投资专家比尔·格罗斯（太平洋投资管理公司）曾经告诉加州大学洛杉矶分校安德森商学院的学生：
\par

我的书房咖啡桌上摆的并不是彼得·林奇的《战胜华尔街》或是我自己的著作，而是历史学家保罗·约翰逊几本有关19世纪和20世纪的历史书……就确定未来而言，没有比历史更好的老师……一本30美元的历史书里隐藏着价值数十亿美元的答案。
\par

\subsubsection*{《枪炮、病菌与钢铁》作者贾雷德·戴蒙德如是说}

“不同民族的历史遵循不同的道路前进，其原因是民族环境的差异，而不是民族自身在生物学上的差异。”
\par

“历史并不‘就是一个又一个讨厌的事实’，就像一个愤世嫉俗者说的那样。的确存在着适用于历史的广泛模式，而寻找对这些模式的解释不但令人陶醉，也是大有裨益的。”
\par

“有一种解释列举了使欧洲人能够屠杀或征服其他民族的直接因素——尤其是欧洲的枪炮、传染病、钢铁工具和工业制成品。这种解释是正确的，然而，这种假设是不全面的，因为它仍然只提供了确定直接原因的一种近似的（初级阶段的）解释。它使人不由得想去寻找终极原因：为什么最后带来枪炮、凶恶的病菌和钢铁的竟是欧洲人，而不是非洲人或印第安人？”
\par
